\chapter{Adaptador FUSE y observabilidad}
\label{cap:fuse}

Este capítulo cierra el loop: una syscall POSIX que el usuario hace
contra el mount point sotFS, ¿cómo se traduce a una regla DPO sobre el
TypeGraph y luego de vuelta en una respuesta? La respuesta vive en
\citaCrate{sotfs-fuse} y conecta todo lo anterior: el grafo del
cap.~\ref{cap:typegraph}, las reglas del cap.~\ref{cap:dpo}, la
persistencia del cap.~\ref{cap:persistencia}, las capabilities del
cap.~\ref{cap:capabilities} y las transacciones del
cap.~\ref{cap:transacciones}.

\section{El callback FUSE: anatomía}

El crate \citaCrate{sotfs-fuse} implementa \texttt{Filesystem} del
crate \texttt{fuser}. Cada syscall mapeable se traduce en un callback
con la siguiente forma:

\begin{algorithm}[H]
\SetAlgoLined
\KwIn{Request $req$, parámetros POSIX, ReplyXxx}
\caption{\textsc{Skeleton de un callback FUSE en sotFS}}\label{alg:fuse-cb}


$\textit{ctx} \gets \textsc{CtxFromReq}(req)$\;
$g.\textsc{SetCapCtx}(\textit{ctx})$\;


$\textit{handles} \gets g.\textsc{Resolve}(\text{paths/inodes del req})$\;


$\textit{result} \gets \textsc{DpoOp}(g, \textit{handles}, \ldots)$\;


\eIf{$\textit{result} = \texttt{Ok}$}{
  $\textit{reply}.\texttt{ok}(\dots)$\;
}{
  $\textit{reply}.\texttt{error}(\textsc{MapErrToErrno}(\textit{result}))$\;
}


\If{\texttt{SOTFS\_PROV\_SIDECAR} configurado}{
  $\textsc{PersistProvLog}(g)$\;
}
\end{algorithm}

\section{Mapeo VFS $\to$ DPO}

\begin{table}[h]
\centering
\renewcommand{\arraystretch}{1.2}
\begin{tabularx}{\linewidth}{l X}
  \toprule
  \textbf{VFS callback} & \textbf{DPO op invocada} \\
  \midrule
  \texttt{create} & \opcreate{} (\citaLine{sotfs-ops/src/lib.rs}{33}) \\
  \texttt{mkdir}  & \opmkdir{} (\citaLine{sotfs-ops/src/lib.rs}{94}) \\
  \texttt{unlink} & \opunlink{} (\citaLine{sotfs-ops/src/lib.rs}{322}) \\
  \texttt{rmdir}  & \oprmdir{} (\citaLine{sotfs-ops/src/lib.rs}{198}) \\
  \texttt{rename} & \oprename{} (\citaLine{sotfs-ops/src/lib.rs}{404}) \\
  \texttt{link}   & \oplink{} (\citaLine{sotfs-ops/src/lib.rs}{272}) \\
  \texttt{read}   & \op{read\_data} (\citaLine{sotfs-ops/src/lib.rs}{704}) \\
  \texttt{write}  & \op{write\_data} (\citaLine{sotfs-ops/src/lib.rs}{670}) \\
  \texttt{setxattr}/\texttt{getxattr}/\texttt{listxattr} & \op{*xattr} \\
  \texttt{symlink}/\texttt{readlink} & \op{symlink}/\op{readlink} \\
  \texttt{getattr}/\texttt{setattr} & lookup directo + \op{chmod}/\op{chown}/\op{truncate} \\
  \bottomrule
\end{tabularx}
\caption{Cada VFS callback de \fuse{} se traduce en una llamada a una
o más funciones DPO de \citaCrate{sotfs-ops}.}
\label{tab:vfs-dpo}
\end{table}

\section{El \texttt{CapContext} desde una syscall}

Recordando del cap.~\ref{cap:capabilities}, cada op corre bajo un
\texttt{CapContext}. La traducción FUSE $\to$ context vive en
\citaLine{sotfs-fuse/src/fs.rs}{59}:

\begin{lstlisting}[language=rust]
fn ctx_from_req(req: &Request<'_>) -> CapContext {
    CapContext::new(None, req.uid() as u64)
}
\end{lstlisting}

Versión actual minimalista: \texttt{cap\_id = None} y
\texttt{domain\_id = uid}. Una extensión natural (no implementada en
v0.2.5) usa el \texttt{pid} o un thread-local con la cap activa
explícita. \citaCommit{61e7faf} cerró el gap más urgente: pre-v0.2.5
el contexto se descartaba al entrar a las DPO ops; ahora se propaga al
\texttt{record\_prov}.

\section{Provenance log: el sidecar JSONL}

\begin{defn}[ProvenanceEntry]
\label{def:prov-entry}
Un \emph{provenance entry} es un registro JSONL con la siguiente
estructura (\citaLine{sotfs-graph/src/provenance.rs}{59}):

\begin{lstlisting}[language=rust]
struct ProvenanceEntry {
    t: u64,          // unix timestamp seconds
    op: ProvOp,      // {Create, Mkdir, Unlink, ...}
    inode: u64,
    cap: Option<u64>,
    domain: u64,
    detail: String,
}
\end{lstlisting}
\end{defn}

\paragraph{Cuándo se escribe.} Si la env var
\texttt{SOTFS\_PROV\_SIDECAR} apunta a un path, cada DPO op exitosa
emite una entry. El sidecar es append-only JSONL: una línea por entry.

\paragraph{El bug v0.2.3 $\to$ v0.2.4.} Pre-v0.2.4 el formato escribía
\texttt{op:Create} (sin comillas) y \texttt{cap:Some(7)}; no era JSON
válido. v0.2.4 cambió a \texttt{serde\_json::to\_string} que produce
JSON spec-compliant: \texttt{"op":"Create"}, \texttt{"cap":7}.
Parsers legacy rompen; consumidores nuevos no.

\section{Graph Hunter: exporter para threat-hunting}

\citaCrate{sotfs-cli} tiene un binario \texttt{sotfs-export-hunter}
con dos modos:

\subsection*{Snapshot mode}

\begin{lstlisting}[language=shellcmd]
sotfs-export-hunter /tmp/test.redb -o hunter.json
\end{lstlisting}

Carga el grafo de \redb, lo convierte a formato \emph{Graph Hunter
temporal multigraph}, y emite un JSON con dos secciones:
\texttt{meta} y \texttt{events}.

\begin{lstlisting}
{
  "meta": {
    "format": "graph-hunter-temporal",
    "version": 1,
    "node_count": 9,
    "edge_count": 10,
    "node_types": ["inode","directory","capability","block","version"],
    "edge_types": ["contains","grants","delegates",
                   "derivesFrom","supersedes","pointsTo","hasXattr"]
  },
  "events": [
    {"t":1778345083, "op":"add_node", "id":"I1",
     "type":"inode", "vtype":"dir", "link_count":1, ...},
    {"t":1778345083, "op":"add_edge", "id":1,
     "src":"D1", "tgt":"I1", "type":"contains", "name":"."}
  ]
}
\end{lstlisting}

\subsection*{Tail mode (streaming)}

\begin{lstlisting}[language=shellcmd]
sotfs-export-hunter --tail /tmp/test.prov.jsonl
sotfs-export-hunter --tail /tmp/test.prov.jsonl --once   # batch
\end{lstlisting}

Tail-f sobre el sidecar JSONL escrito por \texttt{SOTFS\_PROV\_SIDECAR},
convierte cada entry a evento Graph Hunter NDJSON y emite a stdout.
Útil para pipelining a un SIEM:

\begin{lstlisting}[language=shellcmd]
sotfs-export-hunter --tail /tmp/sotfs.prov.jsonl | \
  jq -c '. | select(.cap == 42)'
\end{lstlisting}

\section{Mount: opciones runtime}

\begin{lstlisting}[language=shellcmd]
./target/release/sotfs-fuse /tmp/mnt --db /tmp/test.redb
\end{lstlisting}

Variables de entorno:

\begin{tabularx}{\linewidth}{l X}
  \toprule
  Variable & Efecto \\
  \midrule
  \texttt{SOTFS\_FUSE\_TTL\_MS} & TTL del cache de FUSE (default 1000 ms; 0 desactiva). \\
  \texttt{SOTFS\_FUSE\_ALLOW\_OTHER} & Si set, monta con \texttt{AllowOther + AutoUnmount}; default off (per-UID isolation). \\
  \texttt{SOTFS\_PROV\_SIDECAR} & Path donde escribir el JSONL de provenance. \\
  \bottomrule
\end{tabularx}

\section{Resumen del capítulo}

\begin{itemize}
  \item Cada VFS callback de \fuse{} en \citaCrate{sotfs-fuse} sigue el
    patrón: set cap-context $\to$ resolve $\to$ DPO op $\to$ reply
    $\to$ drain prov-log.
  \item \texttt{ctx\_from\_req} construye el \texttt{CapContext} desde
    el \texttt{Request}; \citaCommit{61e7faf} cerró el plumbing
    hasta el prov-log.
  \item El provenance log es un JSONL sidecar opcional activado por
    \texttt{SOTFS\_PROV\_SIDECAR}. Cada DPO op exitosa emite una entry.
  \item \texttt{sotfs-export-hunter} convierte snapshot o tail-streaming
    a JSON Graph Hunter temporal multigraph, consumible por SIEMs.
\end{itemize}

\section*{Ejercicios}

\begin{ejercicio}[\dificultad{$\star$}]
\label{ej:fuse-1}
Liste las tres variables de entorno que controlan el comportamiento
de \texttt{sotfs-fuse}.
\end{ejercicio}

\begin{ejercicio}[\dificultad{$\star$}]
\label{ej:fuse-2}
¿Qué op DPO invoca el callback \texttt{rename} de FUSE?
\end{ejercicio}

\begin{ejercicio}[\dificultad{$\star\star$}]
\label{ej:fuse-3}
Diseñe un consumidor de provenance que cuente cuántos \opwrite{}
ocurrieron en la última hora por dominio. Pista: jq + grep + sort.
\end{ejercicio}

\begin{ejercicio}[\dificultad{$\star\star$}]
\label{ej:fuse-4}
El \texttt{TTL\_MS} default es 1000 ms. ¿Qué pasa si se setea a 0?
¿Por qué un usuario querría hacerlo (pista: benchmarking)?
\end{ejercicio}

\begin{ejercicio}[\dificultad{$\star\star\star$}]
\label{ej:fuse-5}
Extienda \texttt{ctx\_from\_req} para usar un thread-local de cap
activa explícita, no sólo el \texttt{uid}. ¿Qué hace falta agregar al
runtime para que el FUSE daemon mantenga la map ``thread $\to$ cap''?
\end{ejercicio}
